\documentclass[11pt]{article}

\usepackage[margin=1in]{geometry}
\usepackage{amsmath,amssymb,mathtools}
\usepackage{booktabs,multirow,array}
\usepackage{graphicx}
\usepackage{subcaption}
\usepackage[font=small,labelfont=bf]{caption}
\usepackage{microtype}
\usepackage{xcolor}
\usepackage{enumitem}
\usepackage[numbers,sort&compress]{natbib}
\usepackage{xurl}
\usepackage[colorlinks=true,allcolors=blue!55!black]{hyperref}

\graphicspath{{pinn_solution_benchmark/reports/figures/}}
\setlength{\emergencystretch}{2em}
\setlist{nosep,leftmargin=*}

\newcommand{\ii}{\mathrm{i}}
\newcommand{\Ltwo}{\mathrm{L2RE}}
\newcommand{\sech}{\operatorname{sech}}
\newcommand{\csch}{\operatorname{csch}}
\newcommand{\Description}[1]{} % accessibility description retained in source

\title{\textbf{From Symbolic Audit to Scientific Benchmark:}\\
Verified Travelling Waves and Physics-Informed Neural Networks for a
Fourth-Order Gain--Loss Nonlinear Schr\"odinger Equation}
\author{Authors withheld for review}
\date{September 2, 2026}

\begin{document}
\maketitle

\begin{abstract}
Physics-informed neural networks (PINNs) are frequently validated against exact
solutions, but an incorrect analytic catalogue or residual implementation can
make apparently precise neural predictions scientifically meaningless.  We
therefore couple an exact symbolic audit with a reproducible forward, inverse,
and transfer-learning PINN benchmark for a fourth-order gain--loss nonlinear
Schr\"odinger-type equation containing second- through fourth-order dispersion,
finite gain bandwidth, Kerr and two-photon nonlinearities, and saturable
response.  A proposed $7\times7$ iB-function catalogue is first converted to
polynomial identities and exhaustively solved.  Only 9 of 49 index pairs admit
nontrivial constant-phase solution families; 40 reduce to degenerate closures,
and none of the nine surviving families supports the full physical model.
Every accepted family passes symbolic cancellation and 50-digit evaluation of
the original complex PDE, with relative residuals below $8.4\times10^{-51}$.
Independent sixth-order finite differences with DOP853 reproduce the three
smooth benchmark profiles with best complex relative errors from
$1.6\times10^{-13}$ to $7.5\times10^{-11}$.

The forward benchmark comprises 90 runs over three branches, two
closure-consistent parameter sets, three training strategies, and five seeds.
At the CPU-scaled pilot budget, the constant and dark-kink branches achieve
100\% success for every method under the criterion $\Ltwo<10^{-2}$, whereas
the localized bright $\sech$ branch succeeds only with adaptive loss weighting
(50\%).  Increasing the $\sech$ budget reduces its median error from
$1.6\times10^{-2}$ to $3.5\times10^{-3}$ and raises overall success to 83\%,
showing that the apparent failure is largely budget-dependent.  Inverse
experiments prove that $(g_0,T_2,\alpha)$ are not separately identifiable and
recover instead $c=g_0T_2^2$ and $\ell=(g_0-\alpha)/2$; at 5\% observation
density, their median errors are $1.8$--$6.4\%$ for $c$ and $0.33$--$2.2\%$
for $\ell$ across 0--5\% noise.  A parameter-conditioned PINN fails zero-shot
branch transfer ($\Ltwo\approx0.7$--$1.3$) but becomes competitive with
from-scratch training after fine-tuning on only 1\% of held-out data.  These
results establish a central methodological conclusion: exact-solution
verification, identifiability analysis, independent residual evaluation, and
multi-seed reporting must precede claims about PINN accuracy or
generalization.
\end{abstract}

\noindent\textbf{Keywords:} physics-informed neural networks; nonlinear
Schr\"odinger equation; exact-solution audit; higher-order dispersion;
inverse problems; identifiability; adaptive sampling; scientific benchmarking

%======================================================================
\section{Introduction}
\label{sec:introduction}
%======================================================================

Physics-informed neural networks encode differential equations as soft
constraints in the training objective and thereby combine sparse observations
with a continuous neural representation of the solution.  The modern PINN
formulation demonstrated forward solution and inverse discovery for nonlinear
PDEs, including the nonlinear Schr\"odinger equation (NLSE)
\citep{raissi2019}.  In nonlinear optics, PINNs have since been applied to
standard and higher-order NLSE dynamics, including bright solitons, breathers,
rogue waves, parameter discovery, and parameter-conditioned propagation
\citep{zhou2021,jiang2022,thulasidharan2024}.  Recent work has introduced
conservation penalties \citep{wu2022optical}, multi-domain residual
architectures for interacting waves \citep{chen2025}, and complex-symmetric,
dynamic-weighted networks for high-order NLSE parameter identification
\citep{su2026}.

The analytical component of this study belongs to the
Bogning--Djeumen Tchaho--Kofan\'e methodology (BDKm) developed by Professor
Jean Roger Bogning and collaborators.  Its documented progression begins with
identification and cancellation of hyperbolic-function coefficients
\citep{djeumen2010}, continues through the BDK multi-soliton construction
\citep{djeumen2011}, and is explicitly formulated and named BDKm for
Ginzburg--Landau equations by Bogning et al. \citep{bogning2012}.  A subsequent generalization
extended the construction and its stability analysis \citep{bogning2012general}.
Bogning later systematized the implicit Bogning functions (iB-functions) and
their applications \citep{bogning2019}, including index-based inventories of
probable solution forms \citep{bogning2021}.  Most recently, Ndichia et al.
\citep{ndichia2026} used the same $J_{n,m}$ framework for pump-dependent bright
and dark solitary waves in erbium-doped fibre amplifiers.  These works are the
methodological lineage for the ansatz audited below.

This progress has been accompanied by a sharper understanding of failure.
Differential operators can make the PINN loss severely ill-conditioned even
when the network is expressive enough to represent the target
\citep{krishnapriyan2021,rathore2024}.  Composite data and residual terms can
produce unbalanced gradients \citep{wang2021}, collocation placement can
control accuracy \citep{wu2023}, and global space--time training can violate
the causal structure of evolution equations \citep{wang2024causal}.  Adaptive
residual connections and physics-informed initialization now enable deeper
PINNs to train more reliably \citep{wang2024piratenets}.  These findings imply
that a credible benchmark must report the full seed distribution, distinguish
training loss from error on unseen points, and test whether rankings persist
when the optimization budget changes.

An equally important issue precedes optimization: the reference solution and
the coded residual must be correct.  Exact travelling waves are attractive
ground truth because arbitrary derivatives and parameters can be evaluated
without discretization error.  They are nevertheless unsafe if a symbolic
ansatz is classified by index matching without solving its closure equations,
or if a complex PDE is split incorrectly into real and imaginary parts.  The
candidate catalogue studied here \citep{catalogue2026} proposed 49
iB-function branches for a fourth-order, dissipative, saturable NLSE and a PINN
protocol based on them.  The accompanying reproducibility repository contains
the completed independent audit and numerical experiments needed to test those
claims.

This study makes four contributions.  First, it derives and unit-tests the
correct complex residual split.  Second, it replaces non-rigorous index
matching by an exhaustive polynomial certificate, yielding nine verified
families rather than 49 and explicit rejection reasons for the remainder.
Third, it benchmarks conventional, gradient-balanced, and residual-adaptive
PINNs on the verified solutions, including high-fidelity and stress tests.
Fourth, it evaluates structurally identifiable parameter recovery and
cross-branch transfer.  The work does not claim state-of-the-art accuracy
against unrelated equations or computational budgets; its purpose is to
establish a mathematically audited, failure-aware workflow for this specific
higher-order NLSE.

%======================================================================
\section{Methodology and PINN Architecture}
\label{sec:methods}
%======================================================================

\subsection{Governing equation and identifiable parameterization}

We consider the complex envelope $A(z,t)$ governed by the candidate model
reported in the source catalogue \citep{catalogue2026},
\begin{equation}
\label{eq:pde}
A_z+\frac{\ii}{2}\left(\beta_2+\ii g_0T_2^2\right)A_{tt}
+\frac{\beta_3}{6}A_{ttt}
+\frac{\ii\beta_4}{24}A_{tttt}
+\ii\left(\gamma+\frac{\ii\alpha_2}{2}\right)
\frac{|A|^2A}{1+\Gamma|A|^2}
+\frac{g_0-\alpha}{2}A=0 .
\end{equation}
Here $z$ is propagation distance and $t$ is retarded time;
$\beta_2,\beta_3,\beta_4$ denote second-, third-, and fourth-order dispersion;
$g_0$ is small-signal gain, $T_2$ the gain-medium relaxation time, $\alpha$
linear loss, $\alpha_2$ two-photon-type loss, $\gamma$ Kerr nonlinearity, and
$\Gamma$ nonlinear saturation.  Equation~\eqref{eq:pde} thus has nine physical
coefficients; this terminology follows standard nonlinear-fibre and active-gain
modelling \citep{agrawal2019,ndichia2026}.  Only the combinations
\begin{equation}
\label{eq:reduced-coeffs}
c=g_0T_2^2, \qquad \ell=\frac{g_0-\alpha}{2}
\end{equation}
appear in the residual, so recovering $g_0,T_2,$ and $\alpha$ separately from
a single field is structurally impossible.

Writing $A=u+\ii v$ and
\begin{equation}
Q(u,v)=\frac{u^2+v^2}{1+\Gamma(u^2+v^2)},
\end{equation}
direct expansion of Eq.~\eqref{eq:pde} gives
\begin{align}
F_u={}&u_z-\frac{c}{2}u_{tt}-\frac{\beta_2}{2}v_{tt}
+\frac{\beta_3}{6}u_{ttt}-\frac{\beta_4}{24}v_{tttt}
-\frac{\alpha_2}{2}Qu-\gamma Qv+\ell u,
\label{eq:Fu}\\
F_v={}&v_z-\frac{c}{2}v_{tt}+\frac{\beta_2}{2}u_{tt}
+\frac{\beta_3}{6}v_{ttt}+\frac{\beta_4}{24}u_{tttt}
-\frac{\alpha_2}{2}Qv+\gamma Qu+\ell v.
\label{eq:Fv}
\end{align}
The implementation verifies $F=F_u+\ii F_v$ symbolically on generic real
fields and tests the automatic-differentiation residual against the symbolic
form on random smooth functions.  This corrects the candidate manuscript's
real/imaginary split, which assigned the $\beta_3$ and $\beta_4$ derivatives
to the wrong components.

\subsection{BDKm analytical construction in the source notation}
\label{sec:bdkm-derivation}

The analytical construction follows the Bogning--Djeumen Tchaho--Kofan\'e
methodology: choose a hyperbolic or iB-function prototype, substitute it into
the travelling-wave reduction, use its closed derivative algebra to collect
identical functional terms, and solve the associated coefficient system
\citep{djeumen2010,djeumen2011,bogning2012,bogning2012general}.  Later work
systematized the iB-functions and index-based inventories
\citep{bogning2019,bogning2021}; Ndichia et al. \citep{ndichia2026} use the
same $J_{n,m}$ definition and the pairs $(1,0)$ and $(1,1)$ for bright and dark
waves.  To make the derivation directly comparable with the supplied
\texttt{main.pdf}, this subsection follows its order and notation exactly:
$A_1,\ldots,A_4$ denote third-derivative coefficients and
$C_1,\ldots,C_9$ denote the nine grouped terms.  The reduced variables $c$ and
$\ell$ are not used until the separate verification subsection.

\subsubsection{Travelling-wave transformation}

We first assume that the field depends on the single coordinate
\begin{equation}
\xi=\lambda_1z+\lambda_2t,\qquad A(z,t)=A(\xi),
\label{eq:travelling-coordinate}
\end{equation}
where $\lambda_1$ and $\lambda_2$ are propagation constants.  The chain rule
then gives, without any ansatz-specific simplification,
\begin{equation}
\frac{\partial A}{\partial z}=\lambda_1\frac{dA}{d\xi},\qquad
\frac{\partial^kA}{\partial t^k}=\lambda_2^k\frac{d^kA}{d\xi^k},
\quad k=2,3,4.
\label{eq:chain-transform}
\end{equation}
Substitution into Eq.~\eqref{eq:pde} yields the travelling-wave ordinary
differential equation in the source variables,
\begin{align}
0={}&\lambda_1\frac{dA}{d\xi}
+\frac{\ii}{2}(\beta_2+\ii g_0T_2^2)\lambda_2^2\frac{d^2A}{d\xi^2}
+\frac{\beta_3}{6}\lambda_2^3\frac{d^3A}{d\xi^3}
+\frac{\ii\beta_4}{24}\lambda_2^4\frac{d^4A}{d\xi^4}\nonumber\\
&+\ii\left(\gamma+\frac{\ii\alpha_2}{2}\right)
\frac{|A|^2A}{1+\Gamma|A|^2}+\frac{g_0-\alpha}{2}A,
\label{eq:travelling-ode}
\end{align}
which is Eq.~(8) of the supplied derivation.

\subsubsection{The iB-function ansatz and its derivatives}

The BDKm profile is
\begin{equation}
\label{eq:ansatz}
A(\xi)=aJ_{n,m}(\xi),\qquad
J_{n,m}(\xi)=\frac{\sinh^m\xi}{\cosh^n\xi},
\end{equation}
where $a$ is the wave amplitude.  The defining differentiation rule and the
first derivative are
\begin{equation}
\frac{dJ_{n,m}}{d\xi}=mJ_{n-1,m-1}-nJ_{n+1,m+1},\qquad
\frac{dA}{d\xi}=a[mJ_{n-1,m-1}-nJ_{n+1,m+1}].
\label{eq:J-first}
\end{equation}
Repeated differentiation gives
\begin{align}
\frac{d^2A}{d\xi^2}={}&a\{m(m-1)J_{n-2,m-2}
-[m(n-1)+n(m+1)]J_{n,m}\nonumber\\
&\hspace{33mm}+n(n+1)J_{n+2,m+2}\},
\label{eq:J-second}\\
\frac{d^3A}{d\xi^3}={}&a[A_1J_{n-3,m-3}+A_2J_{n-1,m-1}
+A_3J_{n+1,m+1}+A_4J_{n+3,m+3}],
\label{eq:J-third}
\end{align}
where
\begin{align}
A_1&=m(m-1)(m-2)=m^3-3m^2+2m,\nonumber\\
A_2&=-3m^2n+3m^2-2m,\nonumber\\
A_3&=3mn^2+3n^2+2n,\nonumber\\
A_4&=-n(n+1)(n+2)=-n^3-3n^2-2n.
\label{eq:A-coefficients}
\end{align}
The complete fourth derivative, still in the source convention, is
\begin{align}
\frac{d^4A}{d\xi^4}=a\{&A_1(m-3)J_{n-4,m-4}
+[A_2(m-1)-A_1(n-3)]J_{n-2,m-2}\nonumber\\
&+[A_3(m+1)-A_2(n-1)]J_{n,m}\nonumber\\
&+[A_4(m+3)-A_3(n+1)]J_{n+2,m+2}
-A_4(n+3)J_{n+4,m+4}\}.
\label{eq:J-fourth}
\end{align}
Equations~\eqref{eq:J-first}--\eqref{eq:J-fourth} are precisely the derivative
chain used in the source method; no $T=\tanh\xi$ reparameterization has yet
been introduced.

\subsubsection{Substitution and coefficient grouping}

Using $J_{n,m}^p=J_{pn,pm}$ and dividing Eq.~\eqref{eq:travelling-ode}
by $a$, the source groups the result as
\begin{align}
0={}&C_1J_{n-4,m-4}+C_2J_{n-3,m-3}+C_3J_{n-2,m-2}
+C_4J_{n-1,m-1}+C_5J_{n,m}\nonumber\\
&+C_6J_{n+1,m+1}+C_7J_{n+2,m+2}+C_8J_{n+3,m+3}
+C_9J_{n+4,m+4}\nonumber\\
&+\frac{\ii(\gamma+\ii\alpha_2/2)|a|^2J_{3n,3m}}
{1+\Gamma|a|^2J_{2n,2m}}.
\label{eq:range-rational}
\end{align}
Here the nine coefficients are, in the notation of the supplied derivation,
\begin{align}
C_1={}&\frac{\ii\beta_4}{24}\lambda_2^4m(m-1)(m-2)(m-3),\nonumber\\
C_2={}&\frac{\beta_3}{6}\lambda_2^3m(m-1)(m-2),\nonumber\\
C_3={}&\frac{\ii}{2}(\beta_2+\ii g_0T_2^2)\lambda_2^2m(m-1)
+\frac{\ii\beta_4}{24}\lambda_2^4[A_2(m-1)-A_1(n-3)],\nonumber\\
C_4={}&\lambda_1m+\frac{\beta_3}{6}\lambda_2^3A_2,\nonumber\\
C_5={}&\frac{g_0-\alpha}{2}
-\frac{\ii}{2}(\beta_2+\ii g_0T_2^2)\lambda_2^2[m(n-1)+n(m+1)]\nonumber\\
&+\frac{\ii\beta_4}{24}\lambda_2^4[A_3(m+1)-A_2(n-1)],
\label{eq:C-one-five}\\
C_6={}&-\lambda_1n+\frac{\beta_3}{6}\lambda_2^3A_3,\nonumber\\
C_7={}&\frac{\ii}{2}(\beta_2+\ii g_0T_2^2)\lambda_2^2n(n+1)
+\frac{\ii\beta_4}{24}\lambda_2^4[A_4(m+3)-A_3(n+1)],\nonumber\\
C_8={}&-\frac{\beta_3}{6}\lambda_2^3n(n+1)(n+2),\nonumber\\
C_9={}&\frac{\ii\beta_4}{24}\lambda_2^4n(n+1)(n+2)(n+3).
\label{eq:C-six-nine}
\end{align}

\subsubsection{Fully expanded range equation}

As in \texttt{main.pdf}, multiply Eq.~\eqref{eq:range-rational} by
$1+\Gamma|a|^2J_{2n,2m}$ and distribute every product.  Without summation
shorthand, the complete 19-term range equation is
\begin{align}
0={}&C_1J_{n-4,m-4}+C_2J_{n-3,m-3}+C_3J_{n-2,m-2}
+C_4J_{n-1,m-1}+C_5J_{n,m}\nonumber\\
&+C_6J_{n+1,m+1}+C_7J_{n+2,m+2}+C_8J_{n+3,m+3}
+C_9J_{n+4,m+4}\nonumber\\
&+\Gamma|a|^2C_1J_{3n-4,3m-4}
+\Gamma|a|^2C_2J_{3n-3,3m-3}
+\Gamma|a|^2C_3J_{3n-2,3m-2}\nonumber\\
&+\Gamma|a|^2C_4J_{3n-1,3m-1}
+\Gamma|a|^2C_5J_{3n,3m}
+\Gamma|a|^2C_6J_{3n+1,3m+1}\nonumber\\
&+\Gamma|a|^2C_7J_{3n+2,3m+2}
+\Gamma|a|^2C_8J_{3n+3,3m+3}
+\Gamma|a|^2C_9J_{3n+4,3m+4}\nonumber\\
&+\ii\left(\gamma+\frac{\ii\alpha_2}{2}\right)|a|^2J_{3n,3m}.
\label{eq:range-19}
\end{align}
The first two lines contain the nine unsaturated terms, the next three lines
contain their nine saturation products, and the last line is the nonlinear
term.  This is the exact source equation from which index comparisons are
made.

\subsubsection{Determination of the candidate values of
\texorpdfstring{$n$ and $m$}{n and m}}

Following the supplied derivation, compare a term $J_{n+r,m+r}$ in the first
block of Eq.~\eqref{eq:range-19} with $J_{3n+s,3m+s}$ in the second block,
where
\begin{equation}
r,s\in\{-4,-3,-2,-1,0,1,2,3,4\}.
\end{equation}
Equality of the first indices gives
\begin{equation}
n+r=3n+s\quad\Longrightarrow\quad n=\frac{r-s}{2},
\label{eq:index-n}
\end{equation}
and equality of the second indices gives
\begin{equation}
m+r=3m+s\quad\Longrightarrow\quad m=\frac{r-s}{2}.
\label{eq:index-m}
\end{equation}
For example, $n-4=3n-4$ gives $n=0$, whereas
$n-4=3n-3$ gives $n=-1/2$.  Table~\ref{tab:index-comparisons} reproduces all
81 comparisons used in the source method.

\begin{table}[htbp]
\centering
\caption{Values of $n=(r-s)/2$ obtained from $n+r=3n+s$.  The same table is
obtained for $m$.}
\label{tab:index-comparisons}
\scriptsize
\setlength{\tabcolsep}{4.0pt}
\begin{tabular}{c*{9}{c}}
\toprule
$r\backslash s$ & $-4$ & $-3$ & $-2$ & $-1$ & $0$ & $1$ & $2$ & $3$ & $4$\\
\midrule
$-4$ & $0$ & $-1/2$ & $-1$ & $-3/2$ & $-2$ & $-5/2$ & $-3$ & $-7/2$ & $-4$\\
$-3$ & $1/2$ & $0$ & $-1/2$ & $-1$ & $-3/2$ & $-2$ & $-5/2$ & $-3$ & $-7/2$\\
$-2$ & $1$ & $1/2$ & $0$ & $-1/2$ & $-1$ & $-3/2$ & $-2$ & $-5/2$ & $-3$\\
$-1$ & $3/2$ & $1$ & $1/2$ & $0$ & $-1/2$ & $-1$ & $-3/2$ & $-2$ & $-5/2$\\
$0$  & $2$ & $3/2$ & $1$ & $1/2$ & $0$ & $-1/2$ & $-1$ & $-3/2$ & $-2$\\
$1$  & $5/2$ & $2$ & $3/2$ & $1$ & $1/2$ & $0$ & $-1/2$ & $-1$ & $-3/2$\\
$2$  & $3$ & $5/2$ & $2$ & $3/2$ & $1$ & $1/2$ & $0$ & $-1/2$ & $-1$\\
$3$  & $7/2$ & $3$ & $5/2$ & $2$ & $3/2$ & $1$ & $1/2$ & $0$ & $-1/2$\\
$4$  & $4$ & $7/2$ & $3$ & $5/2$ & $2$ & $3/2$ & $1$ & $1/2$ & $0$\\
\bottomrule
\end{tabular}
\end{table}

Simultaneous matching with the same $(r,s)$ therefore gives
\begin{equation}
n=m=\frac{r-s}{2}.
\label{eq:index-match}
\end{equation}
Although there are 81 comparisons, they generate only the 17 distinct
diagonal candidates
\begin{equation}
(n,m)=(q,q),\quad
q\in\{-4,-7/2,-3,-5/2,-2,-3/2,-1,-1/2,0,
1/2,1,3/2,2,5/2,3,7/2,4\}.
\label{eq:q-set}
\end{equation}

\subsubsection{Probability of appearance in the source construction}

In this part of BDKm, ``probability'' means the frequency with which a pair
appears in the 81 index comparisons.  For $d=r-s=2q$, the number of ordered
pairs with that difference is $N_d=9-|d|$.  Hence
\begin{equation}
N(q)=9-2|q|,\qquad P(q,q)=\frac{N(q)}{81}
=\frac{9-2|q|}{81}.
\label{eq:index-frequency}
\end{equation}
Table~\ref{tab:index-frequency} gives the complete distribution in the same
form as the source, compressed into two side-by-side blocks.

\begin{table}[htbp]
\centering
\caption{Frequency of every diagonal candidate among the 81 comparisons.
This is a combinatorial frequency, not yet an existence probability.}
\label{tab:index-frequency}
\small
\setlength{\tabcolsep}{5pt}
\begin{tabular}{crrc@{\hspace{9mm}}crrc}
\toprule
$q$ & $N$ & $P$ & percent & $q$ & $N$ & $P$ & percent\\
\midrule
$-4$   &1&$1/81$&1.23\% & $1/2$ &8&$8/81$&9.88\%\\
$-7/2$ &2&$2/81$&2.47\% & $1$   &7&$7/81$&8.64\%\\
$-3$   &3&$3/81$&3.70\% & $3/2$ &6&$6/81$&7.41\%\\
$-5/2$ &4&$4/81$&4.94\% & $2$   &5&$5/81$&6.17\%\\
$-2$   &5&$5/81$&6.17\% & $5/2$ &4&$4/81$&4.94\%\\
$-3/2$ &6&$6/81$&7.41\% & $3$   &3&$3/81$&3.70\%\\
$-1$   &7&$7/81$&8.64\% & $7/2$ &2&$2/81$&2.47\%\\
$-1/2$ &8&$8/81$&9.88\% & $4$   &1&$1/81$&1.23\%\\
$0$    &9&$9/81$&11.11\%&       &&&\\
\bottomrule
\end{tabular}
\end{table}
The multiplicities satisfy
$1+2+\cdots+8+9+8+\cdots+2+1=81$, so the tabulated frequencies sum to one.

\subsubsection{Extended and restricted fields of possible solutions}

The supplied derivation next defines the extended field by independently
combining the 17 values in Eq.~\eqref{eq:q-set}, giving $17^2=289$ ordered
pairs.  It then restricts attention to
\begin{equation}
\mathcal I=\{-3/2,-1,-1/2,0,1/2,1,3/2\}
\label{eq:restricted-index-set}
\end{equation}
and forms the $7\times7=49$ ordered pairs in
Table~\ref{tab:restricted-field}.  This is the source document's restricted
``field of possible solutions.''

\begin{table}[htbp]
\centering
\caption{Restricted field obtained by choosing $n$ and $m$ independently from
$\mathcal I$, reproduced in the source convention.}
\label{tab:restricted-field}
\scriptsize
\renewcommand{\arraystretch}{1.25}
\resizebox{\textwidth}{!}{%
\begin{tabular}{c*{7}{c}}
\toprule
$(n,m)$ & $-3/2$ & $-1$ & $-1/2$ & $0$ & $1/2$ & $1$ & $3/2$\\
\midrule
$-3/2$&$(-3/2,-3/2)$&$(-3/2,-1)$&$(-3/2,-1/2)$&$(-3/2,0)$&$(-3/2,1/2)$&$(-3/2,1)$&$(-3/2,3/2)$\\
$-1$&$(-1,-3/2)$&$(-1,-1)$&$(-1,-1/2)$&$(-1,0)$&$(-1,1/2)$&$(-1,1)$&$(-1,3/2)$\\
$-1/2$&$(-1/2,-3/2)$&$(-1/2,-1)$&$(-1/2,-1/2)$&$(-1/2,0)$&$(-1/2,1/2)$&$(-1/2,1)$&$(-1/2,3/2)$\\
$0$&$(0,-3/2)$&$(0,-1)$&$(0,-1/2)$&$(0,0)$&$(0,1/2)$&$(0,1)$&$(0,3/2)$\\
$1/2$&$(1/2,-3/2)$&$(1/2,-1)$&$(1/2,-1/2)$&$(1/2,0)$&$(1/2,1/2)$&$(1/2,1)$&$(1/2,3/2)$\\
$1$&$(1,-3/2)$&$(1,-1)$&$(1,-1/2)$&$(1,0)$&$(1,1/2)$&$(1,1)$&$(1,3/2)$\\
$3/2$&$(3/2,-3/2)$&$(3/2,-1)$&$(3/2,-1/2)$&$(3/2,0)$&$(3/2,1/2)$&$(3/2,1)$&$(3/2,3/2)$\\
\bottomrule
\end{tabular}}
\end{table}

For every entry of this field, the BDKm procedure substitutes $(n,m)$ into
$C_1,\ldots,C_9$ and resolves the remaining algebraic system.  Thus
``possible'' means a profile to be tested by coefficient cancellation; it does
not by itself assert that all physical coefficients can remain active.

\subsection{Closure audit after the source-consistent construction}
\label{sec:closure-audit}

The preceding subsection deliberately reproduces the supplied method.  The
audit is a second layer, introduced only to decide which candidates actually
satisfy the original PDE.  Equations~\eqref{eq:index-n} and
\eqref{eq:index-m} imply diagonal pairs only when the same $(r,s)$ is matched;
forming independent Cartesian products is an exploratory extension rather
than a consequence of that equality.  Furthermore, distinct $J$ indices are
not linearly independent; for example $J_{2,0}=J_{0,0}-J_{2,2}$.  Index
matching is therefore a candidate generator, not a sufficient closure proof.
Conversely, the exact audit accepts the off-diagonal bright branch $(1,0)$,
showing that it is not a necessary condition either.

\subsubsection{Exact closure certificate and admissibility tests}

For this verification layer only, we specialize to the source manuscript's
real constant-phase amplitude $a$, so $|a|^2=a^2$, and use the identifiable
combinations $c=g_0T_2^2$ and $\ell=(g_0-\alpha)/2$ from
Eq.~\eqref{eq:reduced-coeffs}.  Set $T=\tanh\xi$ and $W=\sech\xi$, so that
$J_{n,m}=T^mW^{n-m}$, $W^2=1-T^2$, and
\begin{equation}
D_\xi=(1-T^2)D_T,\qquad L_1=\frac{m}{T}-nT,\qquad
L_{k+1}=(1-T^2)\frac{dL_k}{dT}+L_1L_k,
\label{eq:log-derivative}
\end{equation}
with $J^{(k)}=JL_k$.  Define
\begin{equation}
G(T)=\lambda_1L_1+\frac{\ii}{2}(\beta_2+\ii c)\lambda_2^2L_2
+\frac{\beta_3}{6}\lambda_2^3L_3
+\frac{\ii\beta_4}{24}\lambda_2^4L_4+\ell,
\quad P(T)=J^2.
\end{equation}
Division of Eq.~\eqref{eq:travelling-ode} by nonzero $J$ yields the exact
closure certificate
\begin{equation}
\label{eq:closure-certificate}
[1+\Gamma a^2P(T)]G(T)
+\left(\ii\gamma-\frac{\alpha_2}{2}\right)a^2P(T)\equiv0.
\end{equation}
We split Eq.~\eqref{eq:closure-certificate} into real and imaginary parts.  If
an odd power of $W$ remains, its coefficients are separated in the independent
basis $\{1,W\}$ over rational functions of $T$.  Denominators are cleared only
on an open interval where they are nonzero, and every polynomial coefficient
is set to zero.  A recursive factor-splitting solver enumerates closures in
\begin{equation}
(\lambda_1,\lambda_2,a^2,\beta_2,\beta_3,\beta_4,c,\ell,
\gamma,\alpha_2,\Gamma).
\end{equation}
It rejects $a=0$, $\lambda_2=0$, vacuous equations with at most one active PDE
term, singular saturation denominators, and profiles that are not $C^4$ on
the declared domain.  Half-odd-integer powers receive an explicit principal,
signed-real, or $\xi>0$ convention.

Two familiar profiles illustrate the last algebraic step.  Coefficient
elimination for both forces
$\beta_2=\beta_3=\beta_4=\gamma=\Gamma=\lambda_1=0$, leaving
\begin{equation}
-\frac{c\lambda_2^2}{2}J''+\ell J
-\frac{\alpha_2a^2}{2}J^3=0.
\label{eq:reduced-example}
\end{equation}
For the bright branch $J=J_{1,0}=\sech\xi$,
$J''=J-2J^3$, so Eq.~\eqref{eq:reduced-example} becomes
\begin{equation}
\left(\ell-\frac{c\lambda_2^2}{2}\right)J
+\left(c\lambda_2^2-\frac{\alpha_2a^2}{2}\right)J^3=0,
\end{equation}
and independence of $J$ and $J^3$ gives
$\ell=c\lambda_2^2/2$ and $\alpha_2a^2=2c\lambda_2^2$.  For the dark branch
$J=J_{1,1}=\tanh\xi$, $J''=-2J+2J^3$, producing
\begin{equation}
(\ell+c\lambda_2^2)J
-\left(c\lambda_2^2+\frac{\alpha_2a^2}{2}\right)J^3=0,
\end{equation}
and therefore $\ell=-c\lambda_2^2$ and
$\alpha_2a^2=-2c\lambda_2^2$.  These worked reductions show how a candidate
profile becomes an exact solution family rather than merely a probable index.

Every accepted family passes two final gates: all symbolic coefficients vanish
with free parameters retained, and a rational representative satisfies the
\emph{original complex PDE} at 50-digit precision on a dense $\xi$ grid and
random $(z,t)$ points, with both $\max|F|<10^{-10}$ and
$\|F\|_2/\|A\|_2<10^{-10}$.

\subsection{Independent numerical validation}

Before neural training, each smooth main-group field is propagated by an
independent method of lines: sixth-order central differences in $t$, three
ghost nodes on each side filled from the exact Dirichlet boundary strip, and
DOP853 integration in $z$ \citep{hairer1993}.  Resolutions
$n_t\in\{65,129,257\}$ and tolerances
$10^{-8}$ and $10^{-10}$ form a convergence study.  The kink is never treated
as periodic.  Constant and localized profiles are additionally cross-checked
by a Fourier pseudospectral ETDRK4 method on a three-times-padded temporal
window \citep{kassam2005}.

\subsection{Forward PINN architecture, data, and loss}
\label{sec:forward-method}

The network $\mathcal{N}_{\theta}:(z,t)\mapsto(\hat u,\hat v)$ is a fully
connected multilayer perceptron.  Inputs are affinely scaled to $[-1,1]^2$;
all weights, activations, derivatives, and losses use float64.  The pilot
architecture has four hidden layers of width 40 with $\tanh$ activations,
Xavier-normal weights \citep{glorot2010}, zero biases, and a two-unit linear
output (5,122
trainable parameters).  The smooth activation is essential because the PDE
contains fourth derivatives.  The high-fidelity $\sech$ network uses four
layers of width 48 (7,298 parameters).  No Fourier features are enabled.

For a solution on $[0,Z]\times[t_-,t_+]$, the training set contains
$N_0=256$ exact initial values at $z=0$, $256$ values on each temporal edge
($N_b=512$ total), and $N_f=1024$ Latin-hypercube residual points.  There are
no labelled interior values; the Latin-hypercube construction follows
the original sampling design \citep{mckay1979}.  Validation uses a disjoint $64\times128$ uniform
grid (8,192 field values) and 2,000 newly sampled residual points.  Exact
Dirichlet data are used at both temporal boundaries, including for nonperiodic
kinks.

With $\widehat{\boldsymbol{q}}=(\hat u,\hat v)$ and
$\boldsymbol{q}=(u,v)$, the loss components are
\begin{align}
\mathcal{L}_{0}(\theta)
&=\frac{1}{N_0}\sum_{j=1}^{N_0}
\left\|\widehat{\boldsymbol{q}}_\theta(0,t_j)
-\boldsymbol{q}(0,t_j)\right\|_2^2,
\label{eq:lic}\\
\mathcal{L}_{b}(\theta)
&=\frac{1}{N_b}\sum_{j=1}^{N_b}
\left\|\widehat{\boldsymbol{q}}_\theta(z_j,t_j^{\partial})
-\boldsymbol{q}(z_j,t_j^{\partial})\right\|_2^2,
\quad t_j^{\partial}\in\{t_-,t_+\},
\label{eq:lbc}\\
\mathcal{L}_{f}(\theta)
&=\frac{1}{N_f}\sum_{j=1}^{N_f}
\left(F_u(z_j,t_j;\theta)^2+F_v(z_j,t_j;\theta)^2\right),
\label{eq:lpde}\\
\mathcal{L}(\theta)
&=\lambda_0\mathcal{L}_0+\lambda_b\mathcal{L}_b
+\lambda_f\mathcal{L}_f.
\label{eq:ltotal}
\end{align}
The fixed-weight baseline uses $(\lambda_0,\lambda_b,\lambda_f)=(100,100,1)$.
The gradient-balanced variant updates each data weight every 100 Adam steps by
\begin{equation}
\lambda_k\leftarrow0.9\lambda_k+0.1
\frac{\max|\nabla_\theta\mathcal{L}_f|}
{\operatorname{mean}|\nabla_\theta\mathcal{L}_k|},
\qquad k\in\{0,b\},
\label{eq:gradnorm}
\end{equation}
while retaining unit PDE weight \citep{wang2021}.  The residual-adaptive
distribution (RAD) variant draws $4N_f$ Latin-hypercube candidates and samples
$N_f$ without replacement with
\begin{equation}
p_j=\frac{\varepsilon_j/\bar\varepsilon+1}
{\sum_k(\varepsilon_k/\bar\varepsilon+1)},\qquad
\varepsilon_j=F_{u,j}^2+F_{v,j}^2,
\label{eq:rad}
\end{equation}
at the configured resampling interval \citep{wu2023}.

Pilot optimization uses 600 Adam steps \citep{kingma2015}, initial learning
rate $2\times10^{-3}$
with exponential decay by a factor of 0.1 over the Adam phase, followed by 100
strong-Wolfe L-BFGS iterations (history size 50; gradient and change
tolerances $10^{-12}$ and $10^{-14}$) \citep{liunocedal1989}.  The full matrix contains three
branches, two exact parameter sets, three methods, and seeds 0--4, for 90
runs.  The high-fidelity $\sech$ rerun uses 2,048 residual points, 2,500 Adam
steps, 300 L-BFGS iterations, and seeds 0--2 (18 runs).

The principal metric is complex relative $L^2$ error,
\begin{equation}
\Ltwo=\left[
\frac{\sum_j\left((\hat u_j-u_j)^2+(\hat v_j-v_j)^2\right)}
{\sum_j(u_j^2+v_j^2)+10^{-12}}
\right]^{1/2}.
\label{eq:l2re}
\end{equation}
A run is counted as successful when $\Ltwo<10^{-2}$ and its independently
sampled residual is finite.  Component error, intensity error, spectral error,
maximum pointwise error, initial and boundary RMSE, and residual RMS/max are
also archived.

\subsection{Inverse and conditioned PINNs}

Inverse discovery promotes selected entries of
\begin{equation}
\boldsymbol{\vartheta}=(\beta_2,\beta_3,\beta_4,c,\ell,
\gamma,\alpha_2,\Gamma)
\end{equation}
to trainable scalars.  The loss is
\begin{equation}
\label{eq:inverse-loss}
\mathcal{L}_{\mathrm{inv}}(\theta,\boldsymbol{\vartheta})
=100\,\frac{1}{N_d}\sum_{j=1}^{N_d}
\|\widehat{\boldsymbol{q}}_{\theta}(z_j,t_j)-
\boldsymbol{q}_j^{\mathrm{obs}}\|_2^2
+\mathcal{L}_{f}(\theta,\boldsymbol{\vartheta}).
\end{equation}
The inverse MLP has four width-32 hidden layers; the network and parameters are
jointly optimized for 1,200 Adam steps with learning rates $2\times10^{-3}$
and $5\times10^{-3}$, respectively.  Positive $c$ and $\Gamma$ use softplus
coordinates, all parameters are scaled to order one, and initial guesses are
offset by $\pm50\%$.  The gain/loss sweep uses 1\%, 5\%, and 10\% of a
$48\times96$ grid, corresponding to 46, 230, and 460 observations.  Dispersion,
nonlinear, and joint stages use 230 observations.  Gaussian noise has standard
deviation 0, 1\%, or 5\% of the field-amplitude standard deviation.  Two
branches and three seeds produce 108 completed runs.

Structural identifiability is assessed before recovery by the singular values
of the residual sensitivity matrix at 400 random points.  Both raw
$(g_0,T_2,\alpha)$ and reduced $\boldsymbol{\vartheta}$ coordinates are tested
with a relative rank tolerance $10^{-8}$; sensitivity-matrix rank is a standard
local identifiability diagnostic \citep{iliadis2019}.

For transfer, a conditioned network
$\widehat A_\theta(z,t;\boldsymbol{\eta})$ receives
\begin{equation}
\boldsymbol{\eta}=(n,m,a,\lambda_1,\lambda_2,
\beta_2,\beta_3,\beta_4,c,\ell,\gamma,\alpha_2,\Gamma).
\end{equation}
It has five width-48 $\tanh$ layers and is trained for 800 Adam steps with 512
residual, 96 initial, and 192 boundary values per task.  Interpolation and
extrapolation tasks are generated only by re-solving the exact closure
constraints.  In leave-one-branch-out (LOBO) tests, the omitted branch
contributes no training data.  Fine-tuning and a fresh single-branch baseline
receive the same 250-step budget; fine-tuning sees either 1\% or 5\% of the
held-out interior field.

%======================================================================
\section{Results and Discussion}
\label{sec:results}
%======================================================================

\subsection{The 49-branch claim collapses to nine reduced families}

The exact audit accepts 9 of 49 pairs and rejects 40.  All accepted solutions
are Tier B reduced models; there is no full-model Tier A branch under the
constant-phase ansatz.  Table~\ref{tab:catalogue} reports the complete accepted
catalogue.  The three smooth/bounded branches used in the main benchmark are a
constant, a bright $\sech$, and a dark $\tanh$ kink; the other six are singular
or unbounded stress cases.

\begin{table}[t]
\centering
\caption{All verified constant-phase families.  ``Zero'' counts physical
coefficients forced to zero by closure.  Values are computed at 50-digit
precision from the original complex PDE.}
\label{tab:catalogue}
\small
\begin{tabular}{llllrrr}
\toprule
$(n,m)$ & profile & group & tier & zero & $\max|F|$ & relative RMS \\
\midrule
$(0,0)$   & constant                    & main   & B & 1 & $0$                 & $0$ \\
$(1,0)$   & bright $\sech$              & main   & B & 5 & $1.3\times10^{-51}$ & $5.9\times10^{-52}$ \\
$(1,1)$   & dark $\tanh$ kink           & main   & B & 5 & $2.0\times10^{-51}$ & $6.0\times10^{-52}$ \\
$(-1,-1)$ & $\coth$ (singular)          & stress & B & 5 & $2.1\times10^{-50}$ & $2.4\times10^{-51}$ \\
$(-1,0)$  & $\cosh$ (unbounded)         & stress & B & 2 & $2.1\times10^{-50}$ & $4.8\times10^{-52}$ \\
$(-1,1)$  & $\tfrac12\sinh(2\xi)$      & stress & B & 2 & $3.8\times10^{-48}$ & $2.5\times10^{-51}$ \\
$(0,-1)$  & $\csch$ (singular)          & stress & B & 5 & $2.1\times10^{-50}$ & $3.7\times10^{-51}$ \\
$(0,1)$   & $\sinh$ (unbounded)         & stress & B & 2 & $2.0\times10^{-50}$ & $4.7\times10^{-52}$ \\
$(1,-1)$  & $2\csch(2\xi)$ (singular)  & stress & B & 5 & $4.3\times10^{-50}$ & $8.4\times10^{-51}$ \\
\bottomrule
\end{tabular}
\end{table}

The exact closures of the main branches are
\begin{align}
(0,0):\quad &\gamma=0,\qquad
\ell=\frac{\alpha_2a^2}{2(1+\Gamma a^2)},
\label{eq:closure-constant}\\
(1,0):\quad &\beta_2=\beta_3=\beta_4=\gamma=\Gamma=\lambda_1=0,
\quad \alpha_2a^2=2c\lambda_2^2,
\quad \ell=\tfrac12c\lambda_2^2,
\label{eq:closure-sech}\\
(1,1):\quad &\beta_2=\beta_3=\beta_4=\gamma=\Gamma=\lambda_1=0,
\quad \alpha_2a^2=-2c\lambda_2^2,
\quad \ell=-c\lambda_2^2.
\label{eq:closure-kink}
\end{align}
Thus the bright pulse is a stationary dissipative solution in which spectral
filtering and linear gain are balanced by two-photon absorption; the kink is a
loss-side mirror requiring nonlinear gain.  No nonconstant constant-phase
branch admits saturation or Kerr nonlinearity.  This is a limitation of the
ansatz, not evidence that Eq.~\eqref{eq:pde} lacks such waves.  Related BDKm
analyses also distinguish bright and dark existence regimes through the
dispersion and gain/loss constraints \citep{bogning2012,ndichia2026}.

\begin{figure}[t]
\centering
\begin{subfigure}[t]{0.49\textwidth}
\includegraphics[width=\linewidth]{audit/profile_n1_m0_f0.png}
\caption{Bright $\sech$ branch $(1,0)$.}
\end{subfigure}\hfill
\begin{subfigure}[t]{0.49\textwidth}
\includegraphics[width=\linewidth]{audit/profile_n1_m1_f0.png}
\caption{Dark $\tanh$ branch $(1,1)$.}
\end{subfigure}
\caption{Verified nonconstant main-group profiles at the first archived
representative parameter set.  Each panel shows the real profile, intensity,
and residual magnitude evaluated directly in the complex PDE.}
\Description{Two three-panel plots.  The bright branch is a localized sech
pulse; the dark branch is a monotone tanh kink.  In both, the residual trace is
at the 50-digit arithmetic floor.}
\label{fig:profiles}
\end{figure}

Figure~\ref{fig:surface-heatmaps} shows the same two exact branches over the
full $(z,t)$ validation domain as both three-dimensional surfaces and heat
maps.  The surfaces are intentionally invariant in $z$: closure imposes
$\lambda_1=0$ for these constant-phase branches.  This stationary propagation
is therefore an exact mathematical prediction, not a plotting omission.  The
bright solution is a localized intensity ridge, whereas the dark kink appears
as a persistent central intensity notch.

\begin{figure}[p]
\centering
\includegraphics[width=\textwidth]{summary/bdkm_3d_heatmaps.png}
\caption{Reproducible three-dimensional intensity surfaces (top) and
two-dimensional heat maps (bottom) for the verified bright $\sech$ and dark
$\tanh$ branches.  The plotting program reads $a,\lambda_1,\lambda_2$, the
index pair, and the $\xi$ domain directly from the archived branch catalogue;
the first exact representative set is shown.}
\Description{Four panels show intensity as a function of distance and retarded
time.  The bright branch is a stationary localized ridge and the dark branch
is a stationary central notch.  Lower panels display the same fields as heat
maps.}
\label{fig:surface-heatmaps}
\end{figure}

The source catalogue's matching equations require $n=m$, yet 42 off-diagonal
pairs were added by taking a Cartesian product.  Conversely, the verified
bright branch $(1,0)$ is off-diagonal.  Index matching is therefore neither
necessary nor sufficient.  Repeating the audit with the diagnostic carrier
$A=aJ_{n,m}(\xi)e^{\ii(\kappa z+\omega t)}$ yields 17 families on the same
nine pairs.  It restores conservative travelling bright and dark NLS solitons
and a full-model continuous wave.  The catalogue collapse is consequently
traceable to the claim that a real constant phase is ``without loss of
generality.''

\subsection{Independent solver agreement}

Table~\ref{tab:numerical} summarizes the best finite-difference result and the
independent spectral cross-check.  The localized and kink profiles display
approximately sixth-order spatial convergence before temporal or rounding
error dominates.  The $2.3\times10^{-6}$ ETDRK4 error for the $\sech$ branch
comes from periodic truncation of its nonperiodic tails; the Dirichlet
finite-difference result is $5.8\times10^{-11}$.  Agreement at this level
separates analytic validation from the subsequent PINN evaluation.

\begin{table}[t]
\centering
\caption{Independent numerical validation.  FD--DOP853 entries use
$n_t=257$ and tolerance $10^{-10}$; ETDRK4 uses 256 modes and 400 steps on a
three-times-padded window.}
\label{tab:numerical}
\small
\begin{tabular}{llrr}
\toprule
branch & scheme & final $\Ltwo$ & maximum error \\
\midrule
constant $(0,0)$ & FD--DOP853 & $1.2\times10^{-12}$ & $1.6\times10^{-12}$ \\
                 & ETDRK4     & $8.6\times10^{-18}$ & $8.6\times10^{-18}$ \\
bright $\sech$ $(1,0)$ & FD--DOP853 & $5.8\times10^{-11}$ & $1.2\times10^{-10}$ \\
                 & ETDRK4     & $2.3\times10^{-6}$  & $5.3\times10^{-6}$ \\
dark kink $(1,1)$ & FD--DOP853 & $7.5\times10^{-11}$ & $1.0\times10^{-10}$ \\
\bottomrule
\end{tabular}
\end{table}

\subsection{Forward PINNs: branch difficulty and budget dependence}

All 90 pilot jobs completed without numerical failure.  Table~\ref{tab:forward}
pools the two exact parameter sets for each branch, giving ten runs per
branch--method cell.  The constant branch is easiest, followed by the kink and
then the localized $\sech$.  Fixed weights attain median $\Ltwo$ values
$5.3\times10^{-4}$, $5.3\times10^{-3}$, and $1.9\times10^{-2}$, respectively.
The ordering tracks profile derivative magnitude, especially the fourth
derivative required by Eq.~\eqref{eq:pde}.

\begin{table}[t]
\centering
\caption{Forward pilot matrix pooled over two parameter sets ($n=10$ per
cell).  RMS residuals are evaluated on 2,000 points excluded from training;
the range includes every seed.}
\label{tab:forward}
\footnotesize
\setlength{\tabcolsep}{3.5pt}
\resizebox{\textwidth}{!}{%
\begin{tabular}{llrrrrr}
\toprule
branch & method & median $\Ltwo$ & mean $\pm$ SD & range & success & residual RMS \\
\midrule
constant & fixed
 & $5.3\times10^{-4}$ & $9.7\times10^{-4}\pm7.8\times10^{-4}$
 & $3.4\times10^{-4}$--$2.3\times10^{-3}$ & 100\% & $2.2\times10^{-3}$ \\
$(0,0)$ & adaptive weights
 & $3.8\times10^{-3}$ & $4.1\times10^{-3}\pm3.1\times10^{-3}$
 & $2.6\times10^{-4}$--$8.6\times10^{-3}$ & 100\% & $6.6\times10^{-3}$ \\
 & adaptive sampling
 & $1.3\times10^{-3}$ & $1.6\times10^{-3}\pm1.1\times10^{-3}$
 & $2.6\times10^{-4}$--$3.4\times10^{-3}$ & 100\% & $3.5\times10^{-3}$ \\
\addlinespace
bright $\sech$ & fixed
 & $1.9\times10^{-2}$ & $1.9\times10^{-2}\pm3.3\times10^{-3}$
 & $1.4\times10^{-2}$--$2.4\times10^{-2}$ & 0\% & $2.6\times10^{-2}$ \\
$(1,0)$ & adaptive weights
 & $1.3\times10^{-2}$ & $1.5\times10^{-2}\pm8.4\times10^{-3}$
 & $4.5\times10^{-3}$--$2.9\times10^{-2}$ & 50\% & $1.9\times10^{-2}$ \\
 & adaptive sampling
 & $1.6\times10^{-2}$ & $1.7\times10^{-2}\pm3.6\times10^{-3}$
 & $1.0\times10^{-2}$--$2.4\times10^{-2}$ & 0\% & $2.4\times10^{-2}$ \\
\addlinespace
dark kink & fixed
 & $5.3\times10^{-3}$ & $4.8\times10^{-3}\pm1.1\times10^{-3}$
 & $2.9\times10^{-3}$--$6.1\times10^{-3}$ & 100\% & $1.8\times10^{-2}$ \\
$(1,1)$ & adaptive weights
 & $2.4\times10^{-3}$ & $3.1\times10^{-3}\pm1.7\times10^{-3}$
 & $1.3\times10^{-3}$--$5.9\times10^{-3}$ & 100\% & $9.3\times10^{-3}$ \\
 & adaptive sampling
 & $6.0\times10^{-3}$ & $5.5\times10^{-3}\pm1.5\times10^{-3}$
 & $3.2\times10^{-3}$--$7.4\times10^{-3}$ & 100\% & $1.9\times10^{-2}$ \\
\bottomrule
\end{tabular}
}
\end{table}

\begin{figure}[t]
\centering
\includegraphics[width=0.92\textwidth]{summary/forward_error_distributions.png}
\caption{Complex relative error for every forward seed.  Black bars are
medians, the dashed line is the $10^{-2}$ success threshold, and red diamonds
are the 18 high-fidelity $\sech$ runs.  The display preserves failures and
between-seed dispersion rather than selecting a favourable realization.}
\Description{Log-scale dot plot of every run.  Constant errors lie mostly
below $10^{-3}$, kink errors near several $10^{-3}$, pilot sech errors near or
above $10^{-2}$, and most high-fidelity sech errors below the threshold.}
\label{fig:distributions}
\end{figure}

Adaptive weighting improves the hard $\sech$ and one kink parameterization,
but it worsens the already easy constant solution; RAD is likewise not
uniformly superior.  This agrees with current PINN literature: loss balancing
and adaptive sampling address optimization pathologies but do not define a
problem-independent ranking \citep{wang2021,wu2023,rathore2024}.  In this
benchmark the branch, parameterization, and budget interact materially with
the method.

The high-fidelity rerun confirms that pilot $\sech$ failure is primarily a
budget effect (Table~\ref{tab:hifi}).  Across methods, median $\Ltwo$ falls to
$3.5\times10^{-3}$ and success rises to 83\%.  Fixed weights and RAD reach
100\% success, whereas adaptive weighting falls to 50\%, reversing its pilot
advantage.  For the representative fixed-weight run in
Fig.~\ref{fig:losscurve}, Adam reduces total loss from $3.20\times10^{1}$ to
$1.38\times10^{-3}$ by step 2,499; L-BFGS then lowers it to
$5.44\times10^{-5}$.  Its test $\Ltwo$ is $3.4606\times10^{-3}$ and its
independent residual RMS is $5.5440\times10^{-3}$.

\begin{table}[t]
\centering
\caption{High-fidelity bright-$\sech$ results (two parameter sets, three
seeds, 2,500 Adam plus 300 L-BFGS steps).}
\label{tab:hifi}
\small
\begin{tabular}{lrrrr}
\toprule
method & $n$ & median $\Ltwo$ & range & success \\
\midrule
fixed weights      & 6 & $3.5\times10^{-3}$ & $2.1\times10^{-3}$--$6.5\times10^{-3}$ & 100\% \\
adaptive weights   & 6 & $7.8\times10^{-3}$ & $1.6\times10^{-3}$--$1.6\times10^{-2}$ & 50\% \\
adaptive sampling  & 6 & $2.7\times10^{-3}$ & $2.1\times10^{-3}$--$7.7\times10^{-3}$ & 100\% \\
\midrule
all methods        & 18 & $3.5\times10^{-3}$ & $1.6\times10^{-3}$--$1.6\times10^{-2}$ & 83\% \\
\bottomrule
\end{tabular}
\end{table}

\begin{figure}[t]
\centering
\includegraphics[width=0.72\textwidth]{forward_sech_hifi/losses_forward_sech_hifi__n1_m0_f0__rep0__vanilla__seed2.png}
\caption{Exact archived loss history for the representative high-fidelity
bright-$\sech$ run (parameter set 0, fixed weights, seed 2).  The vertical
change in behaviour at step 2,500 marks the transition from Adam to L-BFGS.
Curves show the unweighted initial, boundary, and PDE components and the
weighted total objective.}
\Description{Semilog loss plot.  Initial and boundary losses rapidly fall
below $10^{-5}$ during Adam; PDE loss decreases more slowly.  L-BFGS produces
a second sharp decrease after step 2,500, ending with data losses below
$2\times10^{-7}$ and PDE loss $3.14\times10^{-5}$.}
\label{fig:losscurve}
\end{figure}

The six stress branches delimit a more difficult regime.  With fixed weights
and three seeds each, median errors range from $1.1\times10^{-2}$ to
$2.9\times10^{-2}$; success is 0--33\%.  Singular branches are trained on
$\xi\in[0.35,4.35]$, and unbounded branches on a finite window.  Their
near-singular fourth derivatives, amplitudes up to approximately 27 in the
selected plotted cases, and active higher-order dispersion explain the loss
of one to two accuracy decades.  These results should not be combined with the
smooth main group when making aggregate claims.

\subsection{Inverse discovery is limited by structural identifiability}

The sensitivity matrix in raw $(g_0,T_2,\alpha)$ coordinates has rank 2 of 3
for both the kink and $\sech$ branches.  The normalized null direction is
$(-0.686,0.243,-0.686)$, confirming Eq.~\eqref{eq:reduced-coeffs}.  Even the
eight-parameter reduced system has numerical rank 7 of 8, with condition
numbers $7.6\times10^{10}$ (kink) and $1.1\times10^{11}$ ($\sech$).  Reporting
nine separately recovered coefficients would therefore overstate what the
field can identify.

Table~\ref{tab:inverse} reports the staged recovery at 5\% density.  The
gain/loss combinations are recovered most reliably, and errors increase
smoothly with noise.  The nonlinear stage estimates nonzero $\alpha_2$ to
1.2--2.4\% and discovers the closure-forced zeros $\gamma$ and $\Gamma$ with
small absolute errors.  Dispersion discovery is weaker, especially for
$\beta_4$, because all three true coefficients are zero on the nonconstant
main branches and the smooth fields provide limited high-order sensitivity.

\begin{table}[t]
\centering
\caption{Median parameter error at 5\% observation density, pooled over two
branches and three seeds per branch.  Entries for true-zero parameters are
absolute errors; all others are relative errors.}
\label{tab:inverse}
\small
\begin{tabular}{llrrrr}
\toprule
stage & parameter & 0\% noise & 1\% noise & 5\% noise & runs \\
\midrule
gain/loss & $c=g_0T_2^2$       & $1.8\times10^{-2}$ & $3.0\times10^{-2}$ & $6.4\times10^{-2}$ & 6 \\
          & $\ell=(g_0-\alpha)/2$ & $3.3\times10^{-3}$ & $7.3\times10^{-3}$ & $2.2\times10^{-2}$ & 6 \\
\addlinespace
nonlinear & $\gamma$ (true 0)  & $1.5\times10^{-3}$ & $7.3\times10^{-4}$ & $2.5\times10^{-3}$ & 6 \\
          & $\Gamma$ (true 0)  & $2.3\times10^{-2}$ & $2.4\times10^{-2}$ & $3.0\times10^{-2}$ & 6 \\
          & $\alpha_2$         & $1.5\times10^{-2}$ & $1.2\times10^{-2}$ & $2.4\times10^{-2}$ & 6 \\
\addlinespace
dispersion & $\beta_2$ (true 0) & $9.1\times10^{-3}$ & $6.9\times10^{-3}$ & $1.8\times10^{-2}$ & 6 \\
           & $\beta_3$ (true 0) & $2.6\times10^{-2}$ & $2.1\times10^{-2}$ & $7.2\times10^{-2}$ & 6 \\
           & $\beta_4$ (true 0) & $4.3\times10^{-2}$ & $6.6\times10^{-2}$ & $7.7\times10^{-2}$ & 6 \\
\bottomrule
\end{tabular}
\end{table}

Increasing observation density from 1\% to 10\% lowers pooled median error
from $4.3\times10^{-2}$ to $1.3\times10^{-2}$ for $c$ and from
$1.8\times10^{-2}$ to $2.5\times10^{-3}$ for $\ell$.  By contrast, joint
four-parameter recovery at 5\% density gives median errors
$1.2\times10^{-1}$ for $c$, $1.6\times10^{-1}$ for $\ell$,
$1.5\times10^{-1}$ for $\alpha_2$, and $1.5\times10^{-2}$ absolute error for
the true-zero $\Gamma$.  The order-of-magnitude degradation of the nonzero
parameters is consistent with the sensitivity condition number rather than a
mere optimizer anomaly.

\begin{figure}[t]
\centering
\includegraphics[width=0.96\textwidth]{summary/inverse_recovery.png}
\caption{Inverse recovery from the 108 archived runs.  Left: median error
against relative observation noise at 5\% density.  Right: gain/loss recovery
against observation density, showing monotone improvement for both identifiable
combinations.}
\Description{Two log-scale line charts.  Errors generally rise with noise;
gain-loss errors decrease from 1 to 10 percent observation density.}
\label{fig:inverse}
\end{figure}

\subsection{Ablations and cross-branch generalization}

Closure-respecting matched ablations are possible only on the constant branch.
Switching saturation from $\Gamma=0$ to 1 changes mean $\Ltwo$ from
$(6.1\pm2.1)\times10^{-4}$ to $(8.8\pm3.0)\times10^{-4}$ ($n=3$ per arm).
Switching $\beta_4$ from 0 to 0.4 changes it from
$(9.2\pm3.5)\times10^{-4}$ to $(1.1\pm0.67)\times10^{-3}$ ($n=3$).  These are
modest effects, and the analytic constant makes the comparison easier than the
nonconstant benchmark.  A conservative ablation cannot be formed: setting
$\ell=\alpha_2=0$ collapses every verified constant-phase family.

The conditioned PINN shows a clear separation between interpolation,
extrapolation, and branch transfer (Table~\ref{tab:generalization}).  Mean
interpolation error on new closure-consistent parameters is
$1.3\times10^{-1}$, while extrapolation error is $6.7\times10^{-1}$.
Zero-shot LOBO error is of order one for every branch, so the conditioning
vector does not induce an unseen qualitative shape.  After seeing only 1\% of
the held-out field, fine-tuning matches or improves on the equal-budget scratch
baseline for all three branches; the benefit is largest for the constant and
$\sech$ cases.  The model has learned reusable features, but not a
branch-complete solution operator.

\begin{table}[t]
\centering
\caption{Conditioned-PINN LOBO transfer.  Values are means over two seeds and
two held-out parameter sets.  Fine-tuning and scratch use the same 250 steps.}
\label{tab:generalization}
\small
\begin{tabular}{lrrrr}
\toprule
held-out branch & zero-shot & fine-tune 1\% & fine-tune 5\% & scratch \\
\midrule
constant $(0,0)$ & $7.1\times10^{-1}$ & $8.3\times10^{-3}$ & $6.2\times10^{-3}$ & $1.1\times10^{-2}$ \\
$\sech$ $(1,0)$  & $1.3\times10^{0}$  & $6.7\times10^{-2}$ & $4.8\times10^{-2}$ & $9.9\times10^{-2}$ \\
kink $(1,1)$     & $1.1\times10^{0}$  & $6.0\times10^{-2}$ & $2.5\times10^{-2}$ & $5.8\times10^{-2}$ \\
\bottomrule
\end{tabular}
\end{table}

\begin{figure}[t]
\centering
\begin{subfigure}[t]{0.54\textwidth}
\includegraphics[width=\linewidth]{summary/generalization_lobo.png}
\caption{Leave-one-branch-out transfer.}
\end{subfigure}\hfill
\begin{subfigure}[t]{0.43\textwidth}
\includegraphics[width=\linewidth]{summary/generalization_interp_extrap.png}
\caption{New parameter values.}
\end{subfigure}
\caption{Parameter-conditioned generalization.  Zero-shot transfer and
extrapolation fail, whereas small-data fine-tuning is competitive at matched
optimization budget.}
\Description{The left grouped bar chart shows order-one zero-shot errors and
much smaller fine-tuned errors.  The right chart shows interpolation errors
below extrapolation errors for all three branches.}
\label{fig:generalization}
\end{figure}

These results refine the literature context.  Parameter inputs can improve
generalization within optical-fiber scenario families \citep{jiang2022}, and
modern architectures can improve optimization and high-order parameter
recovery \citep{wang2024piratenets,su2026}.  Here, however, the dominant
barriers are analytic support and task geometry: only three smooth
constant-phase shapes survive, unseen-branch transfer is discontinuous in
shape, and several coefficients are structurally absent or confounded.  No
network architecture can repair incorrect ground truth or create
identifiability that the governing residual lacks.

\subsection{Scope and limitations}

The experiments were executed on a four-core CPU without a GPU using Python
3.12.3, PyTorch 2.13.0+cpu, NumPy 2.5.2, SciPy 1.18.1, SymPy 1.14.0, and mpmath
1.3.0.  The main numerical claims come from CPU-scaled networks and sample
counts.  A shipped but unexecuted GPU-scale configuration uses six width-64
layers, 10,000 residual points, 384 initial points, 384 points per boundary,
5,000 Adam steps, and 500 L-BFGS iterations.  It is not included in any result
table.  Inverse cells use three of five configured seeds.  The principal
matrices contain 234 completed fits (90 forward-pilot, 18 high-fidelity, 18
stress, and 108 inverse runs), in addition to 12 matched-ablation fits and the
conditioned-PINN transfer fits.  The three smoke-test fits are retained for
software verification but are not pooled into the reported pilot statistics.

The constant-phase ansatz is deliberately narrow.  The extended-phase audit
shows that physically familiar conservative waves return when a carrier is
included, but those families were not mixed into the main benchmark.  The
success threshold $10^{-2}$ is an operational convention, not a theorem.
Moreover, training loss alone is not an error certificate; this is why every
reported run uses disjoint field and residual evaluation.  Finally, metrics
from recent NLSE PINN papers are not directly comparable because equations,
domains, observations, architectures, and budgets differ.  We use that
literature to motivate methods and interpret failure modes, not to claim a
cross-paper leaderboard.

%======================================================================
\section{Conclusion}
\label{sec:conclusion}
%======================================================================

This work turns a proposed exact-solution catalogue into an audited scientific
benchmark.  The symbolic result is decisive: only 9 of 49 constant-phase
iB-function index pairs support nontrivial solutions, all require reduced
physics, and only three are smooth benchmark-grade branches.  The original
PINN residual split is also incorrect; Eqs.~\eqref{eq:Fu}--\eqref{eq:Fv}
provide the tested replacement.  Symbolic cancellation, 50-digit complex-PDE
evaluation, and independent high-order numerical propagation agree before any
neural model is judged.

The PINN results are correspondingly nuanced.  Conventional small MLPs solve
the constant and kink reliably at pilot budget, while the localized $\sech$
requires more optimization and collocation points.  Adaptive weighting helps
at one budget and loses its advantage at another, so single-budget method
rankings are not robust.  Inverse discovery is useful only after reducing to
identifiable coefficient combinations; otherwise the problem asks the network
to distinguish parameters that the PDE itself confounds.  The conditioned
network transfers useful features after minimal fine-tuning but does not
generalize zero-shot to an unseen branch or beyond the trained parameter range.

The practical lesson extends beyond this equation.  A rigorous PINN manuscript
should (i) audit analytic ground truth in the original PDE, (ii) unit-test every
real/imaginary or component residual, (iii) analyze structural
identifiability, (iv) evaluate fields and residuals away from training points,
(v) retain every seed and failed run, and (vi) test whether conclusions survive
a budget change.  Applying this sequence converts attractive loss curves into
claims that are mathematically and computationally defensible.

\section*{Data and Code Availability}

The reproducibility package is contained in
\texttt{pinn\_solution\_benchmark/}.  Exact branch records and rejection
reasons are in \texttt{data/verified\_branches/}; raw per-run JSON files and
checkpoints are in \texttt{results/}; generated reports and figures are in
\texttt{reports/}; configurations are in \texttt{configs/}; and executable
entry points are in \texttt{scripts/}.  The unit-test suite covers symbolic
derivatives, residual identities, automatic differentiation, fourth-order
accuracy, fractional-power conventions, singularity detection, closure
verification, degenerate rejection, metric correctness, data leakage,
determinism, checkpoint/resume, and an end-to-end smoke test.  In particular,
\texttt{scripts/make\_bdkm\_3d\_heatmaps.py} regenerates
Fig.~\ref{fig:surface-heatmaps} from
\texttt{data/verified\_branches/branch\_catalog.json}.

\section*{Acknowledgements}

Author, institutional, funding, and conflict-of-interest information should be
inserted here before submission.

\begin{thebibliography}{99}

\bibitem{raissi2019}
M.~Raissi, P.~Perdikaris, and G.~E. Karniadakis,
``Physics-informed neural networks: A deep learning framework for solving
forward and inverse problems involving nonlinear partial differential
equations,'' \emph{Journal of Computational Physics}, vol.~378, pp.~686--707,
2019. \href{https://doi.org/10.1016/j.jcp.2018.10.045}{doi:10.1016/j.jcp.2018.10.045}.

\bibitem{wang2021}
S.~Wang, Y.~Teng, and P.~Perdikaris,
``Understanding and mitigating gradient flow pathologies in physics-informed
neural networks,'' \emph{SIAM Journal on Scientific Computing}, vol.~43,
no.~5, pp.~A3055--A3081, 2021.
\href{https://doi.org/10.1137/20M1318043}{doi:10.1137/20M1318043}.

\bibitem{krishnapriyan2021}
A.~S. Krishnapriyan, A.~Gholami, S.~Zhe, R.~M. Kirby, and M.~W. Mahoney,
``Characterizing possible failure modes in physics-informed neural networks,''
in \emph{Advances in Neural Information Processing Systems}, vol.~34, 2021.
\url{https://proceedings.neurips.cc/paper/2021/hash/df438e5206f31600e6ae4af72f2725f1-Abstract.html}.

\bibitem{wu2023}
C.~Wu, M.~Zhu, Q.~Tan, Y.~Kartha, and L.~Lu,
``A comprehensive study of non-adaptive and residual-based adaptive sampling
for physics-informed neural networks,'' \emph{Computer Methods in Applied
Mechanics and Engineering}, vol.~403, art.~115671, 2023.
\href{https://doi.org/10.1016/j.cma.2022.115671}{doi:10.1016/j.cma.2022.115671}.

\bibitem{wang2024causal}
S.~Wang, S.~Sankaran, and P.~Perdikaris,
``Respecting causality for training physics-informed neural networks,''
\emph{Computer Methods in Applied Mechanics and Engineering}, vol.~421,
art.~116813, 2024.
\href{https://doi.org/10.1016/j.cma.2024.116813}{doi:10.1016/j.cma.2024.116813}.

\bibitem{rathore2024}
P.~Rathore, W.~Lei, Z.~Frangella, L.~Lu, and M.~Udell,
``Challenges in training PINNs: A loss landscape perspective,'' in
\emph{Proceedings of the 41st International Conference on Machine Learning},
PMLR, vol.~235, pp.~42159--42191, 2024.
\url{https://proceedings.mlr.press/v235/rathore24a.html}.

\bibitem{wang2024piratenets}
S.~Wang, B.~Li, Y.~Chen, and P.~Perdikaris,
``PirateNets: Physics-informed deep learning with residual adaptive networks,''
\emph{Journal of Machine Learning Research}, vol.~25, no.~402, pp.~1--51,
2024. \url{https://jmlr.org/papers/v25/24-0313.html}.

\bibitem{zhou2021}
Z.~Zhou and Z.~Yan,
``Deep learning neural networks for the third-order nonlinear Schr\"odinger
equation: Bright solitons, breathers, and rogue waves,''
\emph{Communications in Theoretical Physics}, vol.~73, art.~105006, 2021.
\href{https://doi.org/10.1088/1572-9494/ac1cd9}{doi:10.1088/1572-9494/ac1cd9}.

\bibitem{jiang2022}
X.~Jiang, D.~Wang, Q.~Fan, M.~Zhang, C.~Lu, and A.~P.~T. Lau,
``Physics-informed neural network for nonlinear dynamics in fiber optics,''
\emph{Laser \& Photonics Reviews}, vol.~16, art.~2100483, 2022.
\href{https://doi.org/10.1002/lpor.202100483}{doi:10.1002/lpor.202100483}.

\bibitem{wu2022optical}
G.~Wu, Y.~Fang, Y.~Wang, Z.~Dai, and C.~Wang,
``Prediction of optical solitons using an improved physics-informed neural
network method with the conservation law constraint,''
\emph{Chaos, Solitons \& Fractals}, vol.~159, art.~112143, 2022.
\href{https://doi.org/10.1016/j.chaos.2022.112143}{doi:10.1016/j.chaos.2022.112143}.

\bibitem{thulasidharan2024}
K.~Thulasidharan, N.~Vishnu Priya, S.~Monisha, and M.~Senthilvelan,
``Predicting positon solutions of a family of nonlinear Schr\"odinger equations
through deep learning algorithm,'' \emph{Physics Letters A}, vol.~511,
art.~129551, 2024.
\href{https://doi.org/10.1016/j.physleta.2024.129551}{doi:10.1016/j.physleta.2024.129551}.

\bibitem{chen2025}
J.~Chen, B.~Li, and M.~W. Yuen,
``Nonlinear interaction multi-domain physics-informed neural networks for the
hybrid rogue wave and breather solutions of the nonlinear Schr\"odinger
equation,'' \emph{Communications in Nonlinear Science and Numerical
Simulation}, vol.~151, art.~109049, 2025.
\href{https://doi.org/10.1016/j.cnsns.2025.109049}{doi:10.1016/j.cnsns.2025.109049}.

\bibitem{su2026}
M.~Su et al.,
``CSD-PINN: A complex-symmetric, dynamic-weighted physics-informed neural
network for parameter identification in high-order nonlinear Schr\"odinger
equations,'' \emph{Nonlinear Dynamics}, vol.~114, art.~185, 2026.
\href{https://doi.org/10.1007/s11071-025-12034-x}{doi:10.1007/s11071-025-12034-x}.

\bibitem{djeumen2010}
C.~T. Djeumen Tchaho, J.~R. Bogning, and T.~C. Kofan\'e,
``Construction of the analytical solitary wave solutions of modified
Kuramoto--Sivashinsky equation by the method of identification of coefficients
of the hyperbolic functions,'' \emph{Far East Journal of Dynamical Systems},
vol.~14, no.~1, pp.~17--34, 2010.
\href{https://doi.org/10.17654/0972111810028}{doi:10.17654/0972111810028}.

\bibitem{djeumen2011}
C.~T. Djeumen Tchaho, J.~R. Bogning, and T.~C. Kofan\'e,
``Multi-soliton solutions of the modified Kuramoto--Sivashinsky equation by
the BDK method,'' \emph{Far East Journal of Dynamical Systems}, vol.~15,
no.~2, pp.~83--98, 2011.
\href{https://www.pphmj.com/journals/articles/786.htm}{Publisher record}.

\bibitem{bogning2012}
J.~R. Bogning, C.~T. Djeumen Tchaho, and T.~C. Kofan\'e,
``Construction of the soliton solutions of the Ginzburg--Landau equations by
the new Bogning--Djeumen Tchaho--Kofan\'e method,'' \emph{Physica Scripta},
vol.~85, no.~2, art.~025013, 2012.
\href{https://doi.org/10.1088/0031-8949/85/02/025013}{doi:10.1088/0031-8949/85/02/025013}.

\bibitem{bogning2012general}
J.~R. Bogning, C.~T. Djeumen Tchaho, and T.~C. Kofan\'e,
``Generalization of the Bogning--Djeumen Tchaho--Kofan\'e method for the
construction of the solitary waves and the survey of the instabilities,''
\emph{Far East Journal of Dynamical Systems}, vol.~20, no.~2,
pp.~101--119, 2012.
\href{https://www.pphmj.com/journals/articles/1033.htm}{Publisher record}.

\bibitem{bogning2019}
J.~R. Bogning, \emph{Mathematics for Physics: The Implicit Bogning Functions
\& Applications}. LAP Lambert Academic Publishing, 2019,
ISBN~978-620-0-23696-8.

\bibitem{bogning2021}
J.~R. Bogning, C.~Jeatsa Dongmo, and C.~Tchawoua,
``The probabilities of obtaining solitary wave and other solutions in the
modified Noguchi power line,'' \emph{Journal of Mathematics Research},
vol.~13, no.~4, pp.~19--29, 2021.
\href{https://doi.org/10.5539/jmr.v13n4p19}{doi:10.5539/jmr.v13n4p19}.

\bibitem{ndichia2026}
L.~O. Ndichia, M.~N. Zambo Abou'ou, M.~Fomekong, and J.~R. Bogning,
``Pump effect on solitary waves in the unsaturated regime of erbium-doped
fiber amplifiers,'' \emph{Journal of Physics Communications}, vol.~10,
no.~2, art.~025005, 2026.
\href{https://doi.org/10.1088/2399-6528/ae4219}{doi:10.1088/2399-6528/ae4219}.

\bibitem{agrawal2019}
G.~P. Agrawal, \emph{Nonlinear Fiber Optics}, 6th~ed. Academic Press,
2019, ISBN~978-0-12-817042-7.

\bibitem{hairer1993}
E.~Hairer, S.~P. N\o rsett, and G.~Wanner,
\emph{Solving Ordinary Differential Equations I: Nonstiff Problems},
2nd~ed. Springer, 1993.

\bibitem{kassam2005}
A.-K. Kassam and L.~N. Trefethen,
``Fourth-order time-stepping for stiff PDEs,''
\emph{SIAM Journal on Scientific Computing}, vol.~26, no.~4,
pp.~1214--1233, 2005.
\href{https://doi.org/10.1137/S1064827502410633}{doi:10.1137/S1064827502410633}.

\bibitem{mckay1979}
M.~D. McKay, R.~J. Beckman, and W.~J. Conover,
``A comparison of three methods for selecting values of input variables in
the analysis of output from a computer code,'' \emph{Technometrics},
vol.~21, no.~2, pp.~239--245, 1979.
\href{https://doi.org/10.1080/00401706.1979.10489755}{doi:10.1080/00401706.1979.10489755}.

\bibitem{glorot2010}
X.~Glorot and Y.~Bengio,
``Understanding the difficulty of training deep feedforward neural networks,''
in \emph{Proceedings of the Thirteenth International Conference on Artificial
Intelligence and Statistics}, PMLR, vol.~9, pp.~249--256, 2010.
\href{https://proceedings.mlr.press/v9/glorot10a.html}{PMLR record}.

\bibitem{kingma2015}
D.~P. Kingma and J.~Ba,
``Adam: A method for stochastic optimization,'' in
\emph{Proceedings of the 3rd International Conference on Learning
Representations}, 2015.
\href{https://arxiv.org/abs/1412.6980}{arXiv:1412.6980}.

\bibitem{liunocedal1989}
D.~C. Liu and J.~Nocedal,
``On the limited memory BFGS method for large scale optimization,''
\emph{Mathematical Programming}, vol.~45, pp.~503--528, 1989.
\href{https://doi.org/10.1007/BF01589116}{doi:10.1007/BF01589116}.

\bibitem{iliadis2019}
A.~Iliadis,
``Structural identifiability and sensitivity,''
\emph{Journal of Pharmacokinetics and Pharmacodynamics}, vol.~46, no.~2,
pp.~127--135, 2019.
\href{https://doi.org/10.1007/s10928-019-09624-9}{doi:10.1007/s10928-019-09624-9}.

\bibitem{catalogue2026}
Anonymous authors,
``Exact travelling-wave solution families of a higher-order gain--loss
nonlinear Schr\"odinger-type equation via the iB-function (BDKm) ansatz: A
combinatorial solution catalog and a benchmark design for physics-informed
machine learning,'' unpublished source manuscript distributed with the
Sim2Sci reproducibility package as \texttt{paper-5.pdf}, 2026.

\end{thebibliography}

\end{document}
